% Fragmento conservado; véase MANIFIESTO_ANTECEDENTES_LECTURA.json.
% Incorporación al final de extension.tex; no contiene una sección principal.
% Procedencia: RESULTADO_RECUPERADO; explicitación sectorial de las hipótesis.
% Fuentes leídas: monografía Doble proyección holográfica (cantera editorial),
% manuscrito/local/01_tesis_recuperada.tex y
% manuscrito/generated/algebra_holonomica.tex, especialmente §§1--7 y 21.
% Definición rectora de número: manuscrito/flattened/main_autosuficiente.tex,
% secciones que comienzan en líneas 22137 y 23000 de la cantera de 775 páginas.
% No se adopta su presentación global como prueba de todas las flechas TPK.

\subsection{Number as a compatible section and value as a readout}
\label{hist:seccion}

The distinction between state and observation makes it possible to specify which object is
prolonged as resolution increases. At depth \(n\), let \(X_n\)
be the set of admissible states produced by APP--TRIT--TPK. Their data
include position, sum and product sheets, residues and quotients, regime,
orientation, carry, phase, route, boundary, memory, and survival. These are not
independent coordinates: preceding arithmetic evaluations and transports
impose their relations. A prolongation preserves those relations
in each antecedent, even as it adds new steps and new memory data.
This is the meaning of the arithmetic of histories developed in
\textup{(provenance preserved in the technical record)}.

Let \(r_{m,n}:X_m\to X_n\), for \(m\geq n\), be the depth
restrictions, with \(r_{n,n}=\operatorname{id}\) and
\(r_{\ell,n}=r_{m,n}\circ r_{\ell,m}\). Each restriction preserves the
prefix and the information already determined in it.

\begin{definition}[Compatible numerical section]
An HMT number is a compatible history of the state system:
\begin{equation}
 x=(x_n)_{n\geq0}\in X_\infty:=\varprojlim_n X_n,
 \qquad r_{m,n}(x_m)=x_n.
 \label{hist:limite-inverso}
\end{equation}
A readout is an additional operation on a specified domain of these
histories. A digit is a finite observation; a scalar value is an
evaluation of the readout. The pair consisting of the section and its readout constitutes
a read presentation, not a redefinition of the section's identity.
\end{definition}

On the Archimedean domain \(X_\infty^{\mathrm{Arch}}\), the readout associates
a closed cylinder with each prefix, preserves their nesting, and makes
their diameter tend to zero. It thus determines
\begin{equation}
 \operatorname{ev}_{\mathbb R}:X_\infty^{\mathrm{Arch}}\longrightarrow
 \mathbb R,\qquad
 \{\operatorname{ev}_{\mathbb R}(x)\}=\bigcap_n I_n(x).
 \label{hist:evaluacion}
\end{equation}
The concrete construction of these cylinders and the determination of their boundaries
are developed in Section~\ref{sec:generacion}. Definition
\eqref{hist:limite-inverso} does not by itself assert that the image of
\eqref{hist:evaluacion} is the entire line: a coverage assertion
requires its own theorem. Nor does equality of values establish, without
an injectivity result, equality of their histories. This separation
makes it possible to speak of precision as depth while preserving the difference
between the arithmetic antecedent and its scalar realization.

\subsection{Composition of histories and memory multiplier}
\label{hist:algebra}

Fix a sector of states at depth \(n\) and a category
\(\mathcal C\) whose arrows are their admissible histories. Composition
is written \(vu=v\circ u\): \(u\) is traversed first and then \(v\).
The elementary steps come from TPK selection, transport, and
updating. A local tritic fiber admits the expression
\(\mathbb R[J_x]/(J_x^2+\tau(x)1_x)\), where
\(\tau(x)\in\{+1,0,-1\}\). Regime and orientation belong to the
state being transported; a transition between regimes is not equivalent to
identifying their quadratic algebras.

The following formulation specifies an algebraic sector of this composition.
For each object \(x\), a unital associative real algebra \(A_x\) is given,
possibly noncommutative. For each \(u:x\to y\), a
unital homomorphism \(\rho_u:A_x\to A_y\) is given. We assume a strict action:
\begin{equation}
 \rho_{\operatorname{id}_x}=\operatorname{id}_{A_x},\qquad
 \rho_v\rho_u=\rho_{vu}.
 \label{hist:accion-estricta}
\end{equation}
For each composable pair \(u:x\to y\), \(v:y\to z\), fix an
invertible central multiplier
\(\Omega(v,u)\in Z(A_z)^\times\). We require the normalization
\(\Omega(\operatorname{id}_y,u)=\Omega(u,\operatorname{id}_x)=1_y\)
and, for each composable triple, the identity
\begin{equation}
 \rho_w\bigl(\Omega(v,u)\bigr)\Omega(w,vu)
   =\Omega(w,v)\Omega(wv,u).
 \label{hist:cociclo}
\end{equation}
These are explicit assumptions on the transports and their coefficients, not a
consequence of calling the dynamics holonomic. When this presentation is applied
to a TPK sector, its \(\rho_u\) and \(\Omega(v,u)\) must be identified.
Transitions expressed through bimodules require that structure to be
preserved and are not replaced by \eqref{hist:accion-estricta}.

The space of finite sums
\[
 \mathcal H(\mathcal C,A,\rho,\Omega)
   =\bigoplus_{u\in\operatorname{Mor}(\mathcal C)} A_{t(u)}T_u
\]
is equipped with the bilinear product
\begin{equation}
 (aT_v)(bT_u)=a\rho_v(b)\Omega(v,u)T_{vu},
 \label{hist:producto}
\end{equation}
if the arrows are composable, and zero otherwise. Here \(t(u)\)
denotes the terminal state. The sum brings together histories with coefficients; the
product concatenates histories while preserving the record multiplier.

\begin{proposition}[Sectorial associativity with a central multiplier]
Under \eqref{hist:accion-estricta} and \eqref{hist:cociclo}, the product
\eqref{hist:producto} is associative. If the set of objects is finite,
its unit is \(\sum_x1_xT_{\operatorname{id}_x}\).
\end{proposition}
\begin{proof}
For \(u:x\to y\), \(v:y\to z\), \(w:z\to t\), let
\(a\in A_y\), \(b\in A_z\), \(c\in A_t\). The two coefficients
of \(T_{wvu}\) are
\begin{align*}
 ((cT_w)(bT_v))(aT_u):\quad&
 c\rho_w(b)\Omega(w,v)\rho_{wv}(a)\Omega(wv,u),\\
 (cT_w)((bT_v)(aT_u)):\quad&
 c\rho_w(b)\rho_w\rho_v(a)\rho_w(\Omega(v,u))\Omega(w,vu).
\end{align*}
The strict action identifies \(\rho_w\rho_v(a)\) with \(\rho_{wv}(a)\).
Centrality of \(\Omega(w,v)\) allows it to be moved to the right of
that factor; \eqref{hist:cociclo} then equates the remaining products.
The coefficients \(a,b,c\) are not permuted. If the triple is not composable,
both parenthesizations are zero by their initial and terminal states.
Bilinearity completes the proof. The local identities and the
normalization of \(\Omega\) verify the indicated unit.
\end{proof}

The calendar carry provides an effective example. In the category
with one object and arrows \(C_9\), take coefficients
\(A=\mathbb R[z,z^{-1}]\), the identity action, and
\[
 \Omega(b,a)=z^{c_0(a,b)},\qquad
 c_0(a,b)=\left\lfloor\frac{a+b}{9}\right\rfloor,
 \quad 0\leq a,b<9.
\]
Both carry sums of a triple equal
\(\lfloor(a+b+c)/9\rfloor\); hence
\eqref{hist:cociclo} holds. The formal variable \(z\) records the winding without
assigning it a numerical value. Subsequently imposing \(z^{12}=1\) recovers the
dodecaphase record of the calendar \eqref{eq:extension108}; retaining
\(z\) without this relation distinguishes successive windings. The example realizes
the memory of this calendar, without identifying it with the entirety of
TPK memory.

\subsection{Double variance of refinement and conservation of unity}
\label{hist:doble-varianza}

Restricting states prolongs observables in the opposite direction.
Consider the function algebras, with pointwise product,
\(D_n=\operatorname{Fun}(X_n,\mathbb R)\). Precomposition defines
\begin{equation}
 r_{m,n}^*:D_n\longrightarrow D_m,
 \qquad r_{m,n}^*f=f\circ r_{m,n}.
 \label{hist:precomposicion}
\end{equation}
These are unital morphisms; they are also injective when the corresponding
restrictions are surjective. The direct limit
\(D_{\mathrm{cyl}}=\varinjlim_n(D_n,r_{n+1,n}^*)\) brings together
observations of finite depth. This algebra is not identified here
with the entire route algebra: prolonging the transport operators also requires
their compatibilities with refinement.

\begin{lemma}[Unity and naturality of evaluation]
If \(e_x\) is the indicator of \(x\in X_n\), then
\begin{equation}
 r_{m,n}^*e_x=\sum_{y:\,r_{m,n}(y)=x}e_y,
 \qquad r_{m,n}^*1_n=1_m,
 \label{hist:unidad}
\end{equation}
and, for every \(f\in D_n\), \(y\in X_m\),
\[
 \langle r_{m,n}^*f,y\rangle=\langle f,r_{m,n}y\rangle.
\]
Each section \(x\in X_\infty\) determines a well-defined evaluation
\(D_{\mathrm{cyl}}\to\mathbb R\), given by \([f]\mapsto f(x_n)\)
for \(f\in D_n\).
\end{lemma}
\begin{proof}
The fibers of \(r_{m,n}\) partition \(X_m\). Their
indicators are orthogonal and sum to \(1_m\), proving
\eqref{hist:unidad}. Naturality is the definition of precomposition.
Finally, \((r_{m,n}^*f)(x_m)=f(x_n)\), so that two compatible
representatives give the same evaluation in the direct limit.
\end{proof}

Global unity is the constant function one over all possibilities;
the individual section selects a compatible history. Preserving the
former does not mean immobilizing the latter. Refinement may increase
the number of extensions and the depth of each history while maintaining
unity and all evaluations of its antecedents.

